% Appendix: the AI work in numbers. The chapter heading is in the main file.
% Every number comes from notes/ai-statistics-data.tex, generated by
% scripts/ai-statistics/make_appendix.py.
\definecolor{aiClaude}{HTML}{4477AA}
\definecolor{aiDeepSeek}{HTML}{EE6677}
\definecolor{aiKimi}{HTML}{CCBB44}
\definecolor{aiGPT}{HTML}{228833}
\definecolor{aiCommits}{HTML}{555555}
\definecolor{aiMilestone}{HTML}{AA3377}

This appendix records how much of the work was done, by which kind of model,
and how often the checks failed. It counts what the project's own records
show. Every number is produced by a short script from files kept in the
repository, and the scripts are kept with it, so a reader can rebuild each
figure and each table. The record covers \aiDays{} days, from \aiFirstDay{}
to \aiLastDay{} (dates in UTC), with \aiCommits{} recorded changes (commits).
The retained reviews-folder census ends on \aiReviewLastDay{}. For \aiLaterFirstDay{}
to \aiLaterLastDay{} the figure also shows a separate count of recovered
review attempts, and the last section lists the reviews of the publication
closeout.

\section*{Who did the work}

The mathematics, the \LaTeX{} source and the bookkeeping were drafted by a
GPT-family model working through Codex. Its work was checked by reviewers from
other model families: Claude (Opus~5.5) at first, and from November~1 DeepSeek
V4 Pro as the main reviewer, with Kimi K3. A GPT model also reviewed the Lean
statements. Appendix~\ref{S:appendix-lean-checks} lists what the Lean companion
checks. A review counts only when the reviewer comes from a different family
from the author of the work it checks.

\begin{center}
\small
\begin{tabular}{@{}lrrr@{}}
  \toprule
  & Sessions or records & Input tokens & Output tokens \\
  \midrule
  Worker: GPT (Codex) & \aiCodexSessions{} sessions & \aiCodexInB{} billion
  & \aiCodexOutM{} million \\
  Claude Code logs & \aiClaudeCalls{} usage records
  & \aiClaudeCacheReadB{} billion (cache) & \aiClaudeOutM{} million \\
  DeepSeek V4 Pro & \aiDeepSeekSessions{} sessions & \aiDeepSeekInM{} +
  \aiDeepSeekCacheM{} million & \aiDeepSeekOutM{} million \\
  \bottomrule
\end{tabular}
\end{center}

Most of the worker's input, \aiCodexCachedPct\%, was text it had already read
and that was served again from a cache; \aiCodexReasonM{} million of its
output tokens were internal reasoning. Counting cached input on both sides,
the worker read about \aiWorkerToDeepSeekRatio{} times as much as the
DeepSeek reviews.
The Claude Code figures include Opus~5.5, Fable~5.2 and synthetic usage
records, covering editing and coordination as well as reviews. They count
logged assistant usage records, not verified API calls. Only the DeepSeek
figures carry a dollar figure,
\$\aiDeepSeekCost{} in US dollars as computed by OpenCode, which is not an
invoice; the other models were used on subscriptions. The worker
machine's logs hold no token counts for Kimi, so these totals leave Kimi out;
the closeout records at the end of this appendix carry Kimi's own reported
counts.
These totals are a snapshot taken at \aiUsageSnapshot{} from those logs.
They cover the sessions and usage records that one machine retained, not all
activity. Only recorded usage fields are summed; usage with no record,
including the abandoned attempt described in the closeout section, is not
included.

\section*{Reviews}

At that historical census, the repository held \aiReviewFolders{} review
folders. Of these,
\aiFoldersWithRun{} contain at least one review run; the others hold
preparation, decisions or diagnostics. A \emph{run} is one review task sent to
one reviewer model, with or without a usable report. There were \aiRuns{}
runs, and \aiFollowups{} further follow-up requests within the same reviewer
sessions.

\begin{center}
\small
\setlength{\tabcolsep}{5pt}
\begin{tabular}{@{}lrrrrrrr@{}}
  \toprule
  Reviewer & Runs & Completed & Not completed & Met bar & FAIL & MEDIUM
  & Hours \\
  \midrule
  \aiReviewTableRows
  \bottomrule
\end{tabular}

\smallskip
\begin{minipage}{0.92\textwidth}
  \footnotesize Hours are summed over the runs whose time was recorded:
  \aiHoursCoverage.
\end{minipage}
\end{center}

\emph{Met the bar} means a completed run graded PASS or PASS WITH NOTES at HIGH
confidence, which is what the project required before a result could be
accepted. Meeting the bar is not the same as acceptance: the final decision was
recorded separately, and only the Claude and GPT reports record it in a form
that can be counted. Of the \aiRunsFail{} FAIL grades, \aiFailGPT{} came from
the GPT reviews of the Lean companion; they reported defects in the Lean
statements or in their match with this book, which then had to be repaired. Runs
that did not complete stopped on usage limits, timeouts or reports in the wrong
format. Figure~\ref{Fig:ai-timeline} shows the work day by day and marks
\aiMilestones{} milestones, from setting up the project to closing the Lean
correspondence gates; their commits are listed in
\path{scripts/ai-statistics/milestones.json}.

The figure covers only this book, which began on \aiInfraBookStart{}. The
system that wrote it is older. The bibliography and the paper library have been
kept since \aiInfraBibStart{}. The collection of \LaTeX\ sources began in
\aiInfraSourcesStart{}, the wiki in its present form in \aiInfraWikiStart{}, and
the library of skills \aiInfraSkillsWhen{}. The book took a month because that
work was already done. The dates and the commits behind them are in
\path{scripts/ai-statistics/infrastructure.json}.

\section*{Approaches that did not work}

As of \aiNogoSnapshot{}, the record holds \aiNogoTotal{} failed approaches
and corrections: \aiNogoMath{} mathematical (\aiNogoRoutes{} approaches that
cannot work and \aiNogoCounter{} counterexamples or corrected claims),
\aiNogoSource{} defects in the source paper, \aiNogoMethod{} failures of a
method of checking, and \aiNogoProcess{} process record.
Chapter~\ref{S:book-no-go} has \aiNogoMath{} numbered mathematical entries
with written arguments at their printed scopes. NG-33 is a broader historical
group of refuted reviewer objections; its printed counterpart proves three
representative cases. The remaining items in that group are not claimed
proved there. In the first snapshot of the record, on November~10,
there were \aiNogoTotalFirst{} entries, and \aiNogoInBookFirst{} of the
mathematical ones were marked as proved in the book, \aiNogoReviewedFirst{} as
reviewed and adjudicated, and \aiNogoRecordedFirst{} as recorded only.
In the current snapshot, \emph{proved in book} denotes a written proof at
the record's stated scope, not scientific acceptance. NG-33 remains
\emph{reviewed and adjudicated} because only part of its historical group is
printed. The earlier review of NG-09 did not meet its acceptance threshold.
None of these records refutes the main theorems.

\begin{figure}[t]
  \centering
  \aiTimelineFigure

  \medskip
  {\scriptsize
  \begin{tabular}{@{}l@{}}
    \aiMilestoneKey
  \end{tabular}\par}
  \caption[The work day by day, with milestones]{Review runs per day by
    reviewer (bars, left scale) and commits per day (line, right scale), with
    the project's \aiMilestones{} milestones (numbered circles). Days are in
    UTC, and each circle sits on the UTC day of the milestone's commit; the
    key gives the date as recorded. Claude returned after November~1 for
    checks of rendered pages. The coloured bars stop on \aiReviewLastDay{},
    where the census of review runs ends; the grey bars count recovered
    review attempts, which are not part of the census. The figure starts with
    this book; the system behind it was built earlier, as the text says.}
  \label{Fig:ai-timeline}
\end{figure}

\section*{What the numbers do not show}

Counts of runs and tokens measure effort, not correctness: a thousand commits do
not prove a theorem, and a passed review is a check, not a proof. Some figures
are incomplete. The worker machine's logs hold no Kimi tokens, and only
\aiRunsKimiTimed{} of the \aiRunsKimi{} Kimi runs have a recorded time; some
DeepSeek runs recorded only their last request, so DeepSeek hours are too low;
and the dates of \aiRunsDatedByFolder{} runs are taken from folder names. The
token totals come from the logs of one machine. The \aiCommits{} commits are
counted in the working repository; the public history has fewer, because commits
that changed only private files are left out of it. A short list of review
folders could not be classified and is kept with the scripts.

The grey bars count \aiLaterAttempts{} review attempts recovered from
retained records for \aiLaterFirstDay{} to \aiLaterLastDay{}. They are
attempts, not completed reviews: some were blocked or returned an unusable
report. Their UTC days use recorded launch, client-session and panel-dispatch
evidence; some are bounded or decoded rather than exact invocation clocks.
Actual model identity is not verified for every attempt, so this series is not
split by reviewer. We show them separately from the \aiRuns{}-run census;
together they still do not establish a complete project-wide inventory. These
are counts of review tasks, not provider requests, and activity with no
retained record is missing. Of these attempts, \aiLaterInLedger{} are
represented in the closeout ledger below.

\section*{The publication closeout reviews of \aiCloseoutMonth}

This section covers only the closeout reviews before publication, in
\aiCloseoutMonth: \aiCloseoutRecords{} records in the experiment's run ledger,
whose snapshot is dated \aiCloseoutSnapshot. The figures come from the dated
correction of that ledger recorded on \aiCloseoutCorrection{}
(\aiCloseoutSource, commit \texttt{\aiCloseoutSourceCommit}, blob
\texttt{\aiCloseoutSourceBlob}). Earlier corrections added verified model
identities and accounting fields and separated one raw report's hash from
that of its parsed copy (C003). The current snapshot preserves all
\aiCloseoutHistoricalRecords{} older
rows unchanged and adds \aiCloseoutRecentRecords{} later cold-review records,
including the author's
bounded decision on the uniform-product review. Reported grades and token
objects remain unchanged. The experiment itself is the way of
doing mathematics described in the Preface, and its primary record is the whole
repository: its history, proofs, failed approaches, human and AI work, sources,
reviews and the Lean companion. The run ledger instruments only this closing
round, and only partly; its own coverage note says what it includes. Its
records are not added to the totals above, and its token counts, reported for
some records only, are not totals for the project. The table counts the outcome
recorded for each record, by the model that actually ran\aiCloseoutRunningNote.
No record accepts the whole book: each acceptance is limited to the scope its
review stated. The preserved older rows retain their historical status: a
row marked as running describes its snapshot date, not a live process now.

\aiCloseoutRecentWords{} later calls add scoped observations in this
November~15 update. A whole-source call returned inconclusive at medium
confidence;
a focused call returned pass with notes at high confidence but did not
complete the required core reconstruction. \aiCloseoutLocalWords{} local
calls returned \aiCloseoutLocalPassWords{} passes and
\aiCloseoutLocalNotesWords{} pass with notes, all at high confidence.
\aiCloseoutLocalAcceptedWords{} were accepted as scoped mathematical evidence
with notes and \aiCloseoutLocalPartialWords{} were partial. For the
uniform-product check, the author accepted a bounded
exception for one automatic native retry, with the incident disclosed. The
original JSON-format failure in one report remains recorded; a
copy in which only the JSON punctuation was corrected is not another review.
These calls do not replace a complete independent review of the book. The
combined critical review of Chapter~12 was not obtained: one attempt
exhausted its context and the next stopped at a safeguard with invalid partial
output. It was never marked passed. A separate high-confidence review of
whether Chapter~12 faithfully states its evidence and limitations does not
replace that missing assessment.
Author-side source comparisons and delivery checks are not additional
independent scientific votes.

The \aiCloseoutRecords{}-record ledger remains a dated snapshot. A separate
evidence addendum links later observations to its original identifiers,
including its missing launch clocks and C011's terminal report. The recovered
inventory contains \aiEvidenceAttempts{} documented review attempts:
\aiEvidenceBaseAttempts{} already represented in this ledger and
\aiEvidenceExtraAttempts{} additional records, with \aiEvidenceFollowups{}
same-session follow-ups listed separately. It is a timeline inventory, not a
new completion, pass-rate or acceptance table; it does not add tokens or costs
to the snapshots above. Its public metadata are in
\aiEvidenceInventoryPath.

\begin{center}
\small
\setlength{\tabcolsep}{5pt}
\aiCloseoutTable
\end{center}

Token counts are given for the \aiCloseoutDenominator{} original review
records, per model, reporting format and accounting scope, in the form in which
the review tool reported them. Only the \aiCloseoutUsageRecords{} records with
reported usage are summed; a record without figures is shown as having no
reported value, and nothing is estimated. Records separately marked as derived
recoveries, parser refusals, preflights or audit addenda do not represent new
provider calls in this ledger. Format continuations count toward usage only
when nonoverlap is verified.\aiCloseoutHeaderNote{} The formats, as recorded
in the ledger:
\begin{itemize}\small
  \setlength{\itemsep}{0pt}
  \aiCloseoutFormatKey
\end{itemize}

Accounting scopes are kept separate:
\begin{itemize}\small
  \setlength{\itemsep}{0pt}
  \aiCloseoutAccountingKey
\end{itemize}
These are sums of known reported components, not actual all-activity totals.
A review record counts a review task, not a provider-stream total. An automatic
retry within one later review record opened two streams of generated output;
it is separate from the recorded repairs to report formats. Usage and cost
for the earlier
abandoned stream are unavailable. The reported final-step total is not the
whole invocation. Actual whole-run usage and cost remain unknown, as does
whether the whole run stayed within its approved token limit.

\begin{center}
\small
\begin{tabular}{@{}>{\raggedright\arraybackslash}p{0.27\textwidth}>{\raggedright\arraybackslash}p{0.5\textwidth}r@{}}
  \toprule
  Model (form) & Reported totals & Reported on \\
  \midrule
  \aiCloseoutTokenRows
  \bottomrule
\end{tabular}
\end{center}

Known reported cost components of the closeout records (not actual
all-activity cost): \aiCloseoutCost.


\subsection*{Later evidence and the author's decision}

Missing cells in the original ledger mean that a value was not verified in
that snapshot, not that no record could ever be recovered. Original launch
records or CLI receipts now supply \aiEvidenceRecordedStarts{} of its
previously missing start clocks:

\begin{center}
\small
\begin{tabular}{@{}lll@{}}
  \toprule
  Record & Recorded clock (UTC) & Kind \\
  \midrule
  \aiEvidenceStartRows
  \bottomrule
\end{tabular}
\end{center}

These are client-runner or CLI receipt clocks, not independently measured
provider API boundaries. They were read from retained records, not calculated
by subtracting a duration from a later clock. The original ledger's nulls
remain unchanged.

For \aiEvidenceClientClocks{} other panel records, the retained clocks record
an OpenCode client session's creation and last update. These do not establish
the review invocation's exact start or finish. For the remaining panel
record, only its UTC day is bounded by the recorded panel creation and the
next role's receipt under the documented sequential workflow. Its exact
start, finish and actual model remain unknown. The evidence addendum keeps
recorded clocks, client clocks, decoded clocks and bounds distinct.

C011 was still running at the original snapshot. Its later receipt records
an invalid report; a punctuation-only recovery returned pass with notes at
medium confidence. The review is terminal and remains unaccepted under the
high-confidence standard. This is later evidence about the same invocation,
not another review or an upgraded grade.

On \aiEvidenceDecisionDate{}, the author chose the policy of an
\emph{honest-standing release with the grades stated plainly}. This decision
changes no grade, supplies no missing review and does not turn scoped
evidence into whole-book scientific acceptance. It selects how the record is
disclosed; it is not a public-release act or the final release accept file.
The dated observation and decision metadata are in \aiEvidenceUpdatesPath.
